% Part: normal-modal-logic
% Chapter: frame-definability
% Section: introduction

\documentclass[../../../include/open-logic-section]{subfiles}

\begin{document}

\olfileid{nml}{frd}{int}

\olsection{Pambuka}

Salah sawijining pitakon kang narik kawigaten para ahli logika
modal yaiku sesambungan antarane relasi aksesibilitas lan kabeneran
!!{formula}s tartamtu ing model kanthi relasi aksesibilitas kasebut.
Upamane, relasi aksesibilitas iku refleksif, yaiku kanggo saben $w
\in W$, $Rww$. Kanthi tembung liya, saben donya bisa diakses saka
awake dhewe. Iki ateges yen $\Box !A$ bener ing sawijining donya~$w$,
$w$ dhewe kalebu donya kang bisa diakses lan mula~$!A$ kudu bener
ing kono. Dadi, yen relasi~$R$ saka~$\mModel{M}$ refleksif, mula
apa wae donya $w$ lan formula $!A$ kang dipilih, $\Box !A
\lif !A$ bener ing kono (kanthi tembung liya, skema $\Box p \lif
p$ lan kabeh instansi substitusine bener ing~$\mModel{M}$).

Nanging kosok baline ora bener. Upamane, yen
$\Box p \lif p$ bener ing~$\mModel{M}$, iki ora ateges $R$ refleksif.
Kita gampang nemokake model~$\mModel{M}$ kang ora refleksif nanging
$\Box p \lif p$ bener ing kabeh donya: pilih model kang mung duwe
siji donya~$w$, kang ora bisa diakses saka awake dhewe, nanging
$w \in V(p)$. Kanthi milih nilai kabeneran $p$ kang cocog, conto iki mung
nuduhake $\Box !A \lif !A$ bener kanggo instansi A minangka p ing model ora refleksif. Cathetan koreksi sumber OLPL-814: dudu kanggo saben A.

Solusine yaiku ngilangi penetapan variabel~$V$ saka
pitungan iki. Yen kita mbutuhake $\Box p \lif p$ bener ing kabeh donya
ing~$\mModel{M}$ tanpa gumantung marang donya endi wae ing~$V(p)$,
mula $R$ kudu refleksif. Sebab ing saben model kang ora refleksif,
ana paling ora siji donya~$w$ kang ora nyukupi $Rww$. Yen kita netepake
$V(p) = W \setminus \{w\}$, mula $p$ bener ing kabeh donya liyane
kajaba~$w$, lan mula uga ing kabeh donya kang bisa diakses saka~$w$
(amarga $w$ mesthi ora bisa diakses saka~$w$, lan $w$ siji-sijine
donya kang~$p$ salah). Nanging, $p$ salah ing $w$, mula $\Box
p \lif p$ salah ing~$w$.

Iki nuduhake yen kita perlu ngenalake notasi kanggo struktur
model tanpa valuasi: iki diarani \emph{bingkai}. Sawijining bingkai
$\mModel{F}$ mung pasangan $\tuple{W, R}$ kang dumadi saka himpunan
donya lan relasi aksesibilitas. Saben model $\tuple{W, R, V}$ banjur,
miturut tetembungan kita, \emph{adhedhasar} bingkai $\tuple{W, R}$.
Kosok baline, sawijining bingkai nemtokake kelas model kang adhedhasar
bingkai iku; lan kelas bingkai nemtokake kelas model kang adhedhasar
salah siji bingkai ing kelas kasebut. Kita bisa nemtokake
$\mModel{F} \Entails !A$, yaiku !!a{formula} kang \emph{valid} ing
sawijining bingkai: $\mSat{M}{!A}$ kanggo kabeh $\mModel{M}$ kang
adhedhasar~$\mModel{F}$.

Kanthi notasi iki, kita bisa netepake relasi korespondensi
antarane !!{formula}s lan kelas bingkai: upamane,
$\mModel{F} \Entails \Box p \lif p$ yen lan mung yen $\mModel{F}$ refleksif.

\end{document}
